% !TEX root = ../main.tex
% Discussed with Roland to change the table. Rationale: methods and analysis are means to reach the contributions, and methods and analysis confuse more than help 
% (from a perspective of someone not so familiar with the paper). Disclosure is not a contribution per se, it's was needed for usenix paper but it's not a contribution of this paper. Also, information ordering is important, so "misconfiguration" for example should be brought more to front. Order the rows chronologically!
\begin{table}[tb]
\centering
\caption{Comparison with related work. ~\cite{landscape} cover ports 389 and 636.}
\scriptsize
\begin{tabular}{l|cccccc|c}
\toprule
& \multicolumn{6}{c|}{\textbf{Contributions \& key insights}} &  \\
 & \rotatebox{90}{\makecell{Hosting and network\\operation of LDAP}} & \rotatebox{90}{\makecell{Use of PKI and\\compromised certs.}} & \rotatebox{90}{\makecell{Misconfigurations and\\sensitive information}} & \rotatebox{90}{\makecell{Recommendations\\for operators}} & \rotatebox{90}{Data set} & \rotatebox{90}{Open-source}
 & \rotatebox{90}{\makecell{Ports 389, 636,\\3268, 3269}} \\
\midrule
Rapid7 (2016)~\cite{project_sonar_rapid7} & \Ocircle & \Ocircle & \Ocircle & \Ccircle & \Ocircle & \Ccircle & \Ccircle \\
Kaspereit (2024)~\cite{landscape} & \Ocircle & \Ocircle & \Ccircle & \Ocircle & \Ocircle & \Ocircle & \Hcircle \\
\midrule
\textbf{This work} & \Ccircle & \Ccircle & \Ocircle & \Ccircle & \Ccircle & \Ccircle & \Ccircle \\
\bottomrule
\end{tabular}
\vspace{-18px}
\label{tab:literature_comparison}
\end{table}





\iffalse
\begin{table*}[ht]
\centering
\caption{Comparison with related work. Half circles are partially done or have been improved in other work.}
\scriptsize
\begin{tabular}{l|cccccc|cccccccccccc}
\toprule
& \multicolumn{6}{c|}{\textbf{Contributions \& key insights}} & \multicolumn{10}{c}{\textbf{Methods \& Analysis}} \\
 & \rotatebox{90}{Hosting of LDAP} & \rotatebox{90}{Disclosure} & \rotatebox{90}{Data set} & \rotatebox{90}{\makecell{Recommendations\\for operators}} & \rotatebox{90}{Open-source} & \rotatebox{90}{Misconfigurations}
 & \rotatebox{90}{Period} & \rotatebox{90}{Ports} & \rotatebox{90}{Identification} & \rotatebox{90}{(start)TLS} & \rotatebox{90}{Extract metadata} & \rotatebox{90}{Network operation} & \rotatebox{90}{Use of PKI} & \rotatebox{90}{Compromised certs.} & \rotatebox{90}{Extract samples} & \rotatebox{90}{Sensitive data} & \rotatebox{90}{LDAP clients} & \rotatebox{90}{UDP scans} \\
\midrule
Kaspereit \etal~\cite{landscape} & \Ocircle & \Ccircle & \Ocircle & \Ocircle & \Ocircle & \Ccircle &2024.01 & 389,636 & \Hcircle & \Ccircle & \Ccircle & \Ocircle & \Ocircle & \Ocircle & \Ccircle & \Ccircle & \Ccircle & \Ocircle \\
Rapid7 blog~\cite{project_sonar_rapid7} & \Ocircle & \Ocircle & \Ocircle & \Ccircle & \Ccircle & \Ocircle & 2016.09 & 389,636,3268,3269 & \Hcircle & \Ocircle & \Ccircle & \Ocircle & \Ocircle & \Ocircle & \Ocircle & \Ocircle & \Ocircle & \Ccircle \\
\midrule
\textbf{This work} & \Ccircle & \Ccircle & \Ccircle & \Ccircle & \Ccircle & \Ocircle & 2024.10--2024.11 & 389,636,3268,3269 & \Ccircle & \Ccircle & \Ccircle & \Ccircle & \Ccircle & \Ccircle & \Ocircle & \Ocircle & \Ocircle & \Ocircle \\
\bottomrule
\end{tabular}
\vspace{-12px}
\label{tab:literature_comparison}
\end{table*}
\fi

\mypara{Related Work}
Empirical studies on the \gls{LDAP} ecosystem are scarce. In \cref{tab:literature_comparison}, we show how our paper is positioned with respect to prior work.
The first work we are aware of is from 2016, when Rapid7's Project Sonar published a blog post~\cite{project_sonar_rapid7} on the presence of \gls{LDAP} servers on the Internet. The blog post did not cover \gls{TLS}, certificates, or hosting.
While the methodology was not precisely described, the company released the \textit{recog} tool, which uses Root DSE entries to identify the version of an \gls{LDAP} server. For our work, we improved \textit{recog} to enhance the software's recognition capabilities. We shared our extensions with the original authors.

The second prior work is our own. In 2024, we introduced \emph{LanDscAPe}, a tool for large-scale \gls{LDAP} security analysis that revealed widespread misconfigurations, leaked passwords, and insecure \gls{TLS} setups~\cite{landscape}, and investigated weaknesses in public-facing \gls{LDAP} servers, focusing on data leaks and misconfigurations.
Here, we expand the scope of our prior work and examine the hosting and deployment practices of \gls{LDAP} servers.
We provide details on the \gls{LDAP} versions, including the product lifecycle, and explore the security practices of \gls{LDAP} servers, including PKI and key management. Our hosting analysis depicts the distribution of \gls{LDAP} servers across different networks. This enables us to investigate whether the \gls{LDAP} ecosystem follows others (like the Web) in terms of centralization. Our analysis includes servers running \gls{GC}

Other kinds of studies have focused on \gls{LDAP} attacks. We do not show these works in \cref{tab:literature_comparison} as they study \gls{LDAP} from a very different perspective and do not measure or characterize the deployment of public-facing servers. For example, Bulusu \etal classified injection attacks and proposed various mitigation techniques~\cite{7210446}. Srinivasa \etal used honeypots to investigate anomalous \gls{LDAP} traffic such as \gls{LDAP} injection, suspicious searches, remote code execution, and brute-force attempts~\cite{9799363}.
When configured over UDP, \gls{LDAP} can be misused for \gls{DDoS} attacks, amplifying traffic by a factor of 56–70, according to Akamai~\cite{WhatCLDAPReflection}. Other works also consider the dangers of in-band upgrades to \gls{TLS}, which include Person-in-the-middle and downgrade attacks~\cite{274731,10.1145/2815675.2815695}.

Prior work focusing on other protocols, Web protocols~\cite{doan2022}, DNS~\cite{Moura2015}, email~\cite{liuWhoGotYour2021}, or \gls{TLS} ~\cite{holz2020} provided evidence for strong Internet centralization. Durumeric \etal investigated the HTTPS ecosystem~\cite{10.1145/2504730.2504755}, and Holz \etal studied protocols for email and chat~\cite{holz2016}.
In contrast, our hosting analysis finds the \gls{LDAP} ecosystem more diversified.
%Other TLS-based protocols have also been studied extensively.